En esta práctica se determinó experimentalmente la relación entre la constante de
Planck y la carga elemental, $h/e$, utilizando diodos emisores de luz (LED) de
diferentes longitudes de onda. Para cada LED se midió el voltaje ánodo--cátodo
correspondiente al umbral de emisión luminosa y, a partir de la longitud de onda
nominal, se calculó la frecuencia de la radiación emitida, permitiendo estimar la relación $h/e$ como $h/e =\left(
\SI{4.13e-15}{} \pm \SI{0.16e-15}{}
\right)\,\si{\volt\second}$ , con un error relativo de $0.22\%$ respecto al valor aceptado.\\


El espectro de emisión de ocho LED se registró mediante un espectrómetro óptico, obteniendo la intensidad luminosa en función de la longitud de onda. Para cada LED, la longitud de onda característica, $\lambda_{\max}$, se identificó como aquella correspondiente al máximo de intensidad medido.\\
